\section*{CHAPTER 2. THEORETICAL BACKGROUND}
\addcontentsline{toc}{section}{\numberline{}CHAPTER 2. THEORETICAL BACKGROUND}

\setcounter{section}{2}
\setcounter{subsection}{0}
\setcounter{figure}{0}
\setcounter{table}{0}
\setcounter{equation}{0}

%============================================================%
\subsection{Overview of Multicarrier Transmission}

Modern broadband wireless systems must combat multipath fading \cite{hanzo2004},
inter-symbol interference (ISI), and limited spectrum resources.
Instead of transmitting all data on a single wideband carrier,
multicarrier systems divide the available bandwidth into many narrowband
subcarriers operating in parallel.

This approach offers two major advantages:

\begin{itemize}
    \item Each subcarrier experiences nearly flat fading.
    \item Equalization complexity is significantly reduced.
\end{itemize}

Orthogonal Frequency Division Multiplexing (OFDM) is the most successful
multicarrier implementation and serves as the foundation for many modern
communication standards.

%============================================================%
\subsection{Orthogonal Frequency Division Multiplexing (OFDM)}

In OFDM, the input bit stream is first grouped and mapped into complex
symbols using modulation schemes such as QPSK or QAM. These symbols are
then assigned to multiple orthogonal subcarriers and converted into a
time-domain waveform using the Inverse Fast Fourier Transform (IFFT).

The transmitted OFDM signal can be expressed as:

\begin{equation}
x[n] = \frac{1}{N}\sum_{k=0}^{N-1} X[k] e^{j2\pi kn/N}
\end{equation}

where:

\begin{itemize}
    \item $X[k]$ is the symbol mapped onto the $k$-th subcarrier
    \item $N$ is the total number of subcarriers
    \item $x[n]$ is the transmitted time-domain sample
\end{itemize}

A Cyclic Prefix (CP) is commonly inserted before transmission to mitigate
inter-symbol interference caused by multipath channels.

The main advantages of OFDM are:

\begin{itemize}
    \item High spectral efficiency
    \item Simple frequency-domain equalization
    \item Robustness against multipath fading
\end{itemize}

However, OFDM suffers from high Peak-to-Average Power Ratio (PAPR).

%============================================================%
\subsection{Single-Carrier Frequency Division Multiple Access (SC-FDMA)}

SC-FDMA can be viewed as a DFT-precoded OFDM system. Before subcarrier
mapping, the modulated data symbols are transformed by a Discrete Fourier
Transform (DFT), which spreads each symbol across the allocated
subcarriers.

Its basic transmitter chain is:

\begin{center}
Bits $\rightarrow$ QAM $\rightarrow$ DFT $\rightarrow$
Subcarrier Mapping $\rightarrow$ IFFT $\rightarrow$ CP
\end{center}

At the receiver, inverse operations are applied:

\begin{center}
Remove CP $\rightarrow$ FFT $\rightarrow$
Subcarrier Demapping $\rightarrow$ IDFT $\rightarrow$ Detection
\end{center}

Because of the DFT spreading operation, SC-FDMA exhibits a waveform with
single-carrier characteristics and lower PAPR than OFDM.

This is the main reason why SC-FDMA was adopted for LTE uplink systems.

%============================================================%
\subsection{Orthogonality of Subcarriers}

The key principle behind OFDM is orthogonality. Subcarriers are spaced so
that they overlap in the frequency domain while remaining separable at the
receiver.

For two different subcarriers $k$ and $m$:

\begin{equation}
\int_{0}^{T}
e^{j2\pi(f_k-f_m)t}dt = 0
\qquad k \neq m
\end{equation}

This property allows efficient spectrum usage without mutual interference
under ideal synchronization conditions.

%============================================================%
\subsection{Peak-to-Average Power Ratio (PAPR)}

PAPR measures how large the signal peaks are relative to its average
power. It is defined as:

\begin{equation}
\mathrm{PAPR}
=
\frac{\max |x[n]|^2}{E\{|x[n]|^2\}}
\end{equation}

High PAPR causes practical power amplifiers to operate inefficiently or
introduce nonlinear distortion.

In OFDM, many independently phased subcarriers may add constructively,
creating large instantaneous peaks. SC-FDMA reduces this effect through
DFT spreading.

%============================================================%
\subsection{Bit Error Rate (BER)}

Bit Error Rate is one of the most common metrics for evaluating digital
communication systems. It is defined as:

\begin{equation}
\mathrm{BER}
=
\frac{\text{Number of erroneous bits}}
{\text{Total transmitted bits}}
\end{equation}

Lower BER indicates more reliable communication performance.

In this report, BER performance under different Signal-to-Noise Ratio
(SNR) conditions will be used to compare OFDM and SC-FDMA.

%============================================================%
\subsection{Comparison Summary}

\begin{table}[H]
\centering
\caption{Comparison between OFDM and SC-FDMA}
\begin{tabularx}{\textwidth}{|>{\raggedright\arraybackslash}X|
>{\centering\arraybackslash}X|
>{\centering\arraybackslash}X|}
\hline
\textbf{Criterion} & \textbf{OFDM} & \textbf{SC-FDMA} \\ \hline
Transmission Type & Multicarrier & Single-carrier-like \\ \hline
PAPR & High & Low \\ \hline
Equalization & Simple & Simple \\ \hline
Complexity & Lower & Slightly Higher \\ \hline
LTE Usage & Downlink & Uplink \\ \hline
\end{tabularx}
\end{table}

%============================================================%
\subsection{Chapter Summary}

This chapter presented the operating principles of OFDM and SC-FDMA,
their mathematical foundations, and the key performance metrics used in
this study. The next chapter will implement both systems in simulation
and compare their BER and PAPR characteristics.